\subsection{关键智能技术的应用}
本节的重点不在于罗列技术名称，而在于说明这些技术在装配过程中的具体作用。结合动力电池包场景，本文选取机器人、人机协作、机器视觉、AI、知识图谱和大模型智能体六类技术进行分析。其中，前四类更直接作用于装配执行和检测过程，后两类更适合作为知识组织与决策辅助工具。

机器人技术首先承担标准化执行者的角色。对于搬运、定位、上料、焊接和重复检查等任务，机器人能够稳定控制节拍并降低操作波动。尤其是焊接和高重复连接任务，由于对路径一致性和姿态控制要求较高，更适合交由机器人完成 \cite{Elhendawy2025IJPEM,Wang2022RCIM}。但动力电池包场景并不意味着完全无人化。由于部件差异、异常返修、线束整理和柔性调整等任务仍具有较强不确定性，人工经验在短期内仍不可替代。

人机协作技术的价值在于重新划分人和机器的边界。对于动力电池包而言，一种更现实的分工方式是：机器人负责重载搬运、重复紧固、固定轨迹检测和标准化焊接，人工负责复杂插接、异常确认、返修和最终判定 \cite{Keshvarparast2024JIM,Gervasi2023IJAMT,Faccio2021IJAMT,Alessio2024JIM,Vysocky2024RCIM,Banziger2023RCIM,Gasparetto2021RCIM,Ji2024SciRep}。这种组合比单纯追求全自动更符合当前制造现场的经济性与可实施性，也更契合智能制造对柔性与安全并重的现实要求。

机器视觉是连接装配与检测的关键桥梁。对于动力电池包场景，视觉系统至少具有三类用途：装配前的定位与缺件检查、过程中的焊缝和装配状态检测，以及终装后的外观与功能关联判定 \cite{Khan2021JRem,Zhu2025JIM,Mayoral2023IJPR,Li2020TIE}。如果缺少视觉检测，许多质量问题便只能依赖人工复检，既影响节拍，也难以保证一致性。

AI 技术在这里不宜狭义理解为“单一神经网络模型”，而应包括缺陷识别、工艺参数优化、异常模式识别和工位级辅助判断等更广义的能力。它更适合作为装配现场的增强工具，而非脱离工艺约束独立决定装配行为 \cite{Nwankpa2021JRem,Mayoral2023IJPR}。

从具体算法结构看，Faster R-CNN 这类两阶段目标检测网络对理解工业视觉系统的工作方式具有代表性。图~\ref{fig:faster-rcnn-structure}清楚展示了共享卷积特征、候选区域生成、RoI 特征提取以及分类和边界框回归等关键环节。其流程通常包括候选区域生成、共享特征提取、目标分类与位置回归等环节，适合承担复杂背景下的缺陷定位和目标级识别任务。对于动力电池包场景，这类网络并不直接替代工艺规则，而是作为视觉感知层的重要组成部分，为后续异常判定、质量追溯和人工复核提供更加结构化的输入 \cite{Ma2024AIR}。

\begin{figure}[H]
  \centering
  \includegraphics[width=0.78\textwidth]{fig_faster_rcnn_structure.png}
  \caption{Faster R-CNN 两阶段检测流程}
  \label{fig:faster-rcnn-structure}
\end{figure}

知识图谱与大模型智能体则承担更高层次的作用。知识图谱适合将产品结构、工艺顺序、质量规则和设备能力组织为统一的知识底座；大模型智能体适合围绕这套知识底座开展问答、异常解释、操作辅助和调度建议。前者解决“知识如何沉淀与表达”的问题，后者解决“知识如何被现场人员快速调用”的问题。在动力电池包装配场景中，这两类技术的价值不在于直接替代机器人或视觉系统，而在于提升现场的理解能力、协同效率和异常处理能力 \cite{Wang2025CJME,Zhang2023JMSYKG,Chen2022AEI,He2023Access,Sgarbossa2024JMSY}。

\begin{figure}[H]
  \centering
  \includegraphics[width=0.76\textwidth]{fig4_fm_hrc_method.png}
  \caption{基础模型支持的人机协作感知与推理框架}
  \label{fig:tech-placeholder}
\end{figure}

图~\ref{fig:tech-placeholder}用于展示人机协作、感知和推理之间的协同关系。图中同时包含语言输入、视觉基础模型感知、机器人 API 调用、执行反馈和推理模块，说明智能装配中的关键技术并不是并列堆放，而是沿着“感知-理解-执行-反馈”链路逐步衔接。这与本文对关键技术协同关系的组织方式一致 \cite{Ji2024SciRep}。
